%! Solving Radical Equations \& Extraneous Solutions — Unit 6, Lesson 6.5 — Homework (Answer Key)
\documentclass[10pt]{article}
\usepackage{algebra2-article}
\usepackage{algebra2-key}   % pulls in -boxes; do NOT also load -boxes
\newcommand{\TallMath}[1]{$\displaystyle #1\rule[-1.4em]{0pt}{3.2em}$}
% In-formula answer slot (\ans is text-mode and must never appear inside $...$).
\newcommand{\mans}[1]{{\color{keyred}\mathbf{#1}}}
\newcommand{\lgrid}[4]{%
  \draw[step=1, linegray!40, very thin] (#1,#3) grid (#2,#4);
  \draw[->, semithick, charcoal] (#1,0)--(#2,0) node[below right, font=\scriptsize]{$x$};
  \draw[->, semithick, charcoal] (0,#3)--(0,#4) node[left, font=\scriptsize]{$y$};}

\begin{document}

\pageheader{Unit 6, Lesson 6.5}{Homework --- Answer Key}
\namedateperiod

\vspace{0.10in}

\begin{notesbox}{Practice}
Isolate, raise, solve, \textbf{check}. Every answer needs a check shown, and ``no real solution'' is
a complete answer when it is the right one.

\begin{enumerate}[leftmargin=1.9em, itemsep=11pt]
  \item \textbf{One radical.} Solve each equation and check.
    \begin{center}
    \renewcommand{\arraystretch}{1.8}
    \begin{tabularx}{\linewidth}{@{}X X X@{}}
    (a) \; $\sqrt{x+6}=4$ \newline $x=\mans{10}$ &
    (b) \; $\sqrt{2x-3}+1=6$ \newline $x=\mans{14}$ &
    (c) \; $3\sqrt{x}-4=8$ \newline $x=\mans{16}$ \\
    (d) \; $\sqrt[3]{x+5}=3$ \newline $x=\mans{22}$ &
    (e) \; $\sqrt[3]{4x}=-2$ \newline $x=\mans{-2}$ &
    (f) \; $\sqrt{x-2}+9=5$ \newline \ans{none} \\
    \end{tabularx}
    \end{center}
    {\small \emph{Work:} (a) $x+6=16$. \; (b) $\sqrt{2x-3}=5$, $2x-3=25$. \;
    (c) $3\sqrt{x}=12$, $\sqrt{x}=4$, $x=16$. \; (d) $x+5=27$. \; (e) $4x=-8$ --- an odd index
    accepts the negative. \; (f) $\sqrt{x-2}=-4$, impossible.}

    Which one of these six needed no algebra at all, and how did you know? \ans{(f) --- once
    isolated it asks a principal square root to equal $-4$, and an even root is never negative.}

  \item \textbf{Variable on both sides.} These are the ones that need sorting. List \emph{every}
        candidate, check each, then give the solution set.
    \begin{center}
    \renewcommand{\arraystretch}{1.7}
    \begin{tabularx}{\linewidth}{@{}l X X X@{}}
    \hline
    \textbf{Equation} & \textbf{Candidates} & \textbf{Extraneous?} & \textbf{Solution set} \\
    \hline
    (a) \; $\sqrt{7-x}=x-1$ & \ans{$3,\ -2$} & \ans{$-2$ is} & \ans{$\{3\}$} \\
    (b) \; $\sqrt{4x+13}=x+4$ & \ans{$-1,\ -3$} & \ans{neither!} & \ans{$\{-3,-1\}$} \\
    (c) \; $x=\sqrt{x+12}$ & \ans{$4,\ -3$} & \ans{$-3$ is} & \ans{$\{4\}$} \\
    \hline
    \end{tabularx}
    \end{center}
    {\small
    \textbf{(a)} $7-x=x^{2}-2x+1 \Rightarrow x^{2}-x-6=0=(x-3)(x+2)$. \;
    Check $3$: $\sqrt4=2=3-1$ \checkmark. \; Check $-2$: $\sqrt9=3$ but $-2-1=-3$. \\
    \textbf{(b)} $4x+13=x^{2}+8x+16 \Rightarrow x^{2}+4x+3=0=(x+1)(x+3)$. \;
    Check $-1$: $\sqrt9=3=-1+4$ \checkmark. \; Check $-3$: $\sqrt1=1=-3+4$ \checkmark. \;
    \emph{Both survive.} \\
    \textbf{(c)} $x^{2}=x+12 \Rightarrow x^{2}-x-12=0=(x-4)(x+3)$. \;
    Check $4$: $\sqrt{16}=4$ \checkmark. \; Check $-3$: $\sqrt9=3$ but $x=-3$.}

    Item (b) is a warning about a habit. Both of its candidates are \emph{negative numbers}, and
    both of them are genuine solutions. So what is it that actually makes a candidate extraneous?
    \ans{Not the candidate's own sign --- it is whether the candidate makes the \emph{non-radical
    side} negative. In (b), $x+4$ is positive at both $-1$ and $-3$, so both check.}

  \item \textbf{Rational exponents.} Raise both sides to the reciprocal power. Watch for a hidden
        $\pm$.
    \begin{center}
    \renewcommand{\arraystretch}{1.8}
    \begin{tabularx}{\linewidth}{@{}X X X X@{}}
    (a) \; $x^{3/2}=8$ \newline $x=\mans{4}$ &
    (b) \; $x^{2/3}=25$ \newline $x=\mans{\pm 125}$ &
    (c) \; $(x-1)^{1/3}=4$ \newline $x=\mans{65}$ &
    (d) \; $x^{4/3}=16$ \newline $x=\mans{\pm 8}$ \\
    \end{tabularx}
    \end{center}
    {\small \emph{Work:} (a) $x=8^{2/3}=\left(\sqrt[3]8\right)^{2}=4$. \;
    (b) $x=\pm 25^{3/2}=\pm\left(\sqrt{25}\right)^{3}=\pm125$; check $(-125)^{2/3}=(-5)^{2}=25$
    \checkmark. \; (c) cube: $x-1=64$. \;
    (d) $x=\pm 16^{3/4}=\pm\left(\sqrt[4]{16}\right)^{3}=\pm8$; check $(-8)^{4/3}=(-2)^{4}=16$
    \checkmark.}

    Two of these have \emph{two} solutions. What do those two have in common? \ans{An \textbf{even
    numerator} --- $x^{2/3}$ and $x^{4/3}$ are squares and fourth powers in disguise, so they lose
    the sign of $x$ and both $\pm$ values work.}

  \item \textbf{Read the answers off a graph.} The curve below is $y=\sqrt{x+2}$.
    \begin{center}
    \begin{tikzpicture}[scale=0.42]
      \lgrid{-4}{15}{-2}{5}
      \begin{scope}
        \clip (-4,-2) rectangle (15,5);
        \draw[very thick, forest, domain=-2:15, samples=110, smooth] plot (\x,{sqrt((\x)+2)});
        \draw[very thick, navy] (-4,3)--(15,3);
        \draw[very thick, goldacc] (-4,1)--(15,1);
        \draw[very thick, redacc] (-4,-1)--(15,-1);
      \end{scope}
      \fill[forest] (-2,0) circle (4pt);
      \fill[charcoal] (7,3) circle (3pt);
      \fill[charcoal] (-1,1) circle (3pt);
      \node[navy, font=\scriptsize, right] at (12.4,3.5) {$y=3$};
      \node[goldacc, font=\scriptsize, right] at (12.4,1.5) {$y=1$};
      \node[redacc, font=\scriptsize, right] at (12.4,-1.5) {$y=-1$};
    \end{tikzpicture}
    \end{center}
    \begin{enumerate}[label=(\alph*), leftmargin=1.9em, itemsep=4pt]
      \item $\sqrt{x+2}=3$ has \ans{one} solution(s): $x=\mans{7}$
      \item $\sqrt{x+2}=1$ has \ans{one} solution(s): $x=\mans{-1}$
      \item $\sqrt{x+2}=-1$ has \ans{no} solution(s), because \ans{the line $y=-1$ lies entirely
            below the curve --- a principal square root is never negative, so they never meet}
      \item Solve (c) algebraically anyway --- square both sides and see what number comes out.
            $x=\mans{-1}$. \; What is that number called, and what does the graph say about it?
            \ans{It is \textbf{extraneous}. Squaring turned $y=-1$ into the same equation as
            $y=1$, so the algebra returned part (b)'s answer --- the graph shows $(-1,1)$ on the
            curve, but nothing at $(-1,-1)$.}
    \end{enumerate}

  \item \textbf{Error analysis.} A classmate turns in this work.
    \begin{center}
    \begin{tcolorbox}[colback=white, colframe=charcoal!45, boxrule=0.7pt, arc=1.5mm,
        width=0.62\linewidth, left=3mm, right=3mm, top=1.5mm, bottom=1.5mm]
    \centering\small
    \renewcommand{\arraystretch}{1.3}
    \begin{tabular}{@{}l@{}}
    $\sqrt{x}+2=5$ \\
    $x+4=25$ \\
    $x=21$ \\
    \end{tabular}
    \end{tcolorbox}
    \end{center}
    \vspace{-4pt}
    (a) Check $x=21$ in the original equation. Does it work? \ans{No --- $\sqrt{21}+2\approx6.58$,
    not $5$.} \\[3pt]
    (b) Which \emph{step} of the method did they skip, and what should line 2 have been?
    \ans{Step 1, isolating. They squared a \emph{sum}: $\left(\sqrt{x}+2\right)^{2}=x+4\sqrt{x}+4$,
    not $x+4$. Line 2 should have been $\sqrt{x}=3$.} \\[3pt]
    (c) Solve it correctly. $x=\mans{9}$

  \item \textbf{Explain.} For a \emph{square root} equation the check can reject an answer, but for
        a \emph{cube root} equation the raising step can never produce one to reject. Explain both
        halves of that sentence, using the words \textbf{index} and \textbf{sign}.
        \ansline{An even \textbf{index} gives the principal root, which is never negative, and
        squaring sends $a$ and $-a$ to the same number --- so squaring loses the \textbf{sign}
        information and the squared equation can be true where the original is false. That is where
        an extraneous candidate comes from, and only the check catches it.\\[3pt]
        An odd \textbf{index} keeps the \textbf{sign} of its radicand, and every real number has
        exactly one real cube root, so $A^{3}=B^{3}$ forces $A=B$. Cubing is reversible: the cubed
        equation has exactly the same solutions as the original and cannot gain extras. (You still
        check --- for arithmetic, not for phantoms.)}
\end{enumerate}
\end{notesbox}

\begin{extensionbox}
\textbf{Reading the skid marks.} Crash investigators estimate how fast a car was going from the
length of the skid marks it left. If a car skids $d$ feet on dry asphalt before stopping, its speed
when the brakes locked was about
\[
  s(d)=\sqrt{24d} \ \text{ miles per hour.}
\]
In Lesson 6.0 you were given $d$ and found $s$. An investigator has the opposite problem: they know
the speed they are testing and need the skid length --- which means solving a radical equation.
\begin{enumerate}[label=(\alph*), leftmargin=1.8em, itemsep=5pt]
  \item A driver claims they were doing the $30$ mph speed limit. Solve $30=\sqrt{24d}$ to find how
        long the skid marks would be at that speed.
  \item Now solve for a speed of $60$ mph. Compare the two skid lengths: the speed doubled, but what
        happened to the distance? Which feature of the graph of $s(d)=\sqrt{24d}$ predicted that
        before you did any algebra?
  \item The actual skid marks at the scene are $96$ feet long. How fast was the car going? Was the
        driver telling the truth?
  \item Solving $s=\sqrt{24d}$ for $d$ in general gives $d=\dfrac{s^{2}}{24}$ --- and this is exactly
        the formula you derived in the Lesson 6.0 homework. You obtained it by squaring both sides,
        which is the step that usually manufactures extraneous solutions. Explain why there is no
        risk of one here.
  \item What does $s(0)=0$ say about the situation, and why does the context rule out negative
        values of $d$ even though the algebra would happily accept them?
\end{enumerate}
\ansline{(a) $30=\sqrt{24d} \Rightarrow 900=24d \Rightarrow d=37.5$ feet.}
\ansline{(b) $3600=24d \Rightarrow d=150$ feet. Doubling the speed \emph{quadrupled} the skid ---
$37.5\to150$. The square root graph flattens as it goes right, so a large change in $d$ buys only a
small change in $s$; run that backwards and a modest gain in speed demands a huge gain in distance.}
\ansline{(c) $s=\sqrt{24(96)}=\sqrt{2304}=48$ mph. The driver was not telling the truth --- $48$ mph
is well over the $30$ mph limit.}
\ansline{(d) Because both sides are non-negative here: $s$ is a speed and $\sqrt{24d}$ is a
principal root, so neither can be the negative twin of the other. Extraneous solutions come from
squaring away a \emph{sign disagreement}, and there is no sign disagreement available.}
\ansline{(e) $s(0)=0$: no skid marks means the car was not moving when the brakes locked. A negative
$d$ is meaningless --- you cannot skid a negative distance --- and the radicand agrees, since
$24d\ge0$ already forces $d\ge0$.}
\end{extensionbox}

\begin{spiralbox}
\textbf{Coming next --- Lesson 6.6: Composition of Functions.} Today you undid a radical by raising
both sides to a power --- one operation cancelling another. That idea is about to get a name and a
notation. In 6.6 you will feed one function's output straight into another function,
$(f\circ g)(x)=f(g(x))$, and find out that the order matters enormously. Then in 6.7 composition
becomes the \emph{test}: two functions are inverses exactly when composing them in either order
hands you back the $x$ you started with --- which will finally explain the reflection over $y=x$
you first noticed all the way back in Lesson 6.0.
\end{spiralbox}

\begin{teachernote}
\small \textbf{Item 2(b) is the most valuable problem on the page.} Both candidates are negative and
both are correct. It is the only reliable cure for ``reject the negative one,'' which is what most
students actually learn if every extraneous root they ever meet happens to be negative. If you grade
one item closely, grade the follow-up sentence to 2(b).
\begin{itemize}[leftmargin=1.3em, nosep]
  \item \textbf{Item 4(d)} is subtle and worth going over: squaring $\sqrt{x+2}=-1$ produces
        \emph{part (b)'s} answer, because squaring cannot tell the line $y=-1$ from the line $y=1$.
        The graph makes the phantom visible.
  \item \textbf{Item 3(b) and 3(d)} will be missed by students who drop the $\pm$. Tie it to
        $x^{2}=25$, which nobody answers with a single number.
  \item \textbf{The extension} closes a loop opened in 6.0, where students derived $d=s^{2}/24$ from
        the same model. Part (d) is the interpretation item for \textbf{A2.EI.5b} --- an honest
        example of squaring being completely safe, which is worth having on record after a lesson
        spent being suspicious of it.
\end{itemize}
\end{teachernote}

\end{document}
